
% https://digital.runeft.ru/6_2020/mtrele

%\needspace{4\baselineskip}
\color{uniblue}{\subsection{Контроль коммутационного ресурса выключателя (КРВ)}}
\color{black}

\begin{enumerate}[label=\arabic{section}.\arabic{subsection}.\arabic*, labelsep=4pt, leftmargin=0pt, itemindent=57pt]

\item
Контроль коммутационного ресурса выключателя предназначен для контроля состояния выключателя в процессе его эксплуатации и определения плановых сроков проведения регламентных работ.

\item
В устройстве контроль коммутационного ресурса выключателя представлен функцией «КРВ».

\item
Функция вычисляет коммутационный (КРВ) и механический ресурсы (МРВ) выключателя и выдает сигнал об их превышении. Коммутационный ресурс проверяется по каждой фазе отдельно. При снижении $\textsl{КРВ}_{(k)}$ по одной из фаз ниже уставки «КРВсраб» функция выдает сигнал о превышении коммутационного ресурса выключателя («Превышен ресурс В»). Счетчик механического ресурса выполняет расчет на каждом цикле включения отключения выключателя. Если количество операций счетчика МРВ превысит уставку «МРВпасп», то функция выдает сигнал о превышении механического ресурса выключателя («Превышен МРВ»).
\item
Анализ механического ресурса выполняется в блоке «МРВ». При вводе функции в работу блок «МРВ» присваивает параметру $\textsl{МРВ}_{(k)}$  значение уставки «Нач\_знач\_МРВ» и сохраняет его в независимую память устройства в качестве параметра \text{$\textsl{МРВ}_{(k-1)}$}. В процессе дальнейшей работы устройства контролируется выполнение цикла включения-отключения (В-О) выключателя:
\begin{itemize}
\item если цикла В-О не произошло, то $\textsl{МРВ}_{(k)}$ = \text{$\textsl{МРВ}_{(k-1)}$} и значение в памяти остается неизменным;
\item если цикл В-О производился, то $\textsl{МРВ}_{(k)}$ = \text{$\textsl{МРВ}_{(k-1)}$} + 1. Значение в памяти должно обновиться на данную величину.
\end{itemize}

\item
Выполнение цикла В-О фиксируется при появлении сигнала «В: Отключен», кратковременно появляющийся (в течении выдержки времени таймера «T2») после переключения блок-контактов выключателя (переход состояния выключателя из положения «Включен» в положение «Отключен»).

\item
Анализ коммутационного ресурса выполняется в блоке «КРВ». При вводе функции в работу, уставке «Нач\_знач\_КРВ» по умолчанию присваивается значение равное 100\%. После чего параметру $\textsl{КРВ}_{(k)}$ в блоке «КРВ» присваивается значение «Нач\_знач\_КРВ» и сохраняется в независимую память устройства в качестве параметра \text{$\textsl{КРВ}_{(k-1)}$}.

\item
Старт расчета $\textsl{КРВ}_{(k)}$ осуществляется при одновременном наличии следующих условий:
\begin{itemize}
\item КРВ введен в работу;
\item наличие сигнала на выходе импульсного таймера T2 (уставка «Tмакс\_В») (действие на электромагниты отключения выключателя);
\item наличие сигнала «В: Отключен».
\end{itemize}

\item
Расчет $\textsl{КРВ}_{(k)}$ осуществляется по максимальному значению тока $I_{\textsl{макс\_выч}}$, полученному на интервале времени работы импульсного таймера T2 (уставка «Tмакс\_В»).
Формула расчета КРВ различна в зависимости от величины тока $I_{\textsl{макс\_выч}}$:

\item
Если 0 < $I_{\textsl{макс\_выч}}$ $\leq$ $I_{\textsl{ном\_В\_пасп}}$, то «КРВ» рассчитывается по формуле:

% ФОРМУЛА №1
\begin{equation*}
\textsl{КРВ}_{(k)} = \text{$\textsl{КРВ}_{(k-1)}$} - \frac{\textsl{100}}{\text{$\textsl{МРВ}_\textsl{пасп}$}\times (\frac{\text{$\textsl{КРВ}_\textsl{пасп\_Iном}$}}{\text{$\textsl{МРВ}_\textsl{пасп}$}})^{\frac{I_{\textsl{макс\_выч}}}{I_{\textsl{ном\_В\_пасп}}}}},
\end{equation*}

Если выполняются неравенства $I_{\textsl{ном\_В\_пасп}}$ $<$  $I_{\textsl{макс\_выч}}$ $\leq$  $I_{\textsl{ном\_откл\_В\_пасп}}$, то «КРВ» рассчитывается по формуле:

\begin{equation*}
\textsl{КРВ}_{(k)} = \text{$\textsl{КРВ}_{(k-1)}$} - \frac{\textsl{100}}{\text{$\textsl{КРВ}_\textsl{пасп\_Iном}$}\times (\frac{\text{$\textsl{КРВ}_\textsl{пасп\_Iном}$}}{\text{$\textsl{КРВ}_\textsl{пасп\_Iном\_откл}$}})^{\frac{ln(\frac{I_{\textsl{ном\_откл\_В\_пасп}}}{I_{\textsl{макс\_выч}}})}{ln(\frac{I_{\textsl{ном\_откл\_В\_пасп}}}{I_{\textsl{ном\_В\_пасп}}})}}},
\end{equation*}

\begin{eqexpl}[25mm] % 55мм только здесь - здесь специфичное описание переменных, обычно 25mm
\item{$\textsl{КРВ}_{(k)}$}значение ресурса выключателя на текущем цикле работы функции;
\item{$\text{$\textsl{КРВ}_{(k-1)}$}$}значение ресурса выключателя на предыдущем цикле работы функции;
\item{$\text{$\textsl{МРВ}_\textsl{пасп}$}$}допустимый механический ресурс выключателя;
\item{$\text{$\textsl{КРВ}_\textsl{пасп\_Iном}$}$}допустимое количество циклов при номинальном токе выключателя;
\item{$\text{$\textsl{КРВ}_\textsl{пасп\_Iном\_откл}$}$}допустимое количество циклов при номинальном токе отключения выключателя;
\item{$I_{\textsl{ном\_В\_пасп}}$}номинальный ток выключателя;
\item{$I_{\textsl{ном\_откл\_В\_пасп}}$}номинальный ток отключения выключателя.

%\item{\begin{minipage}{0.35\textwidth}
%$\text{$\textsl{МРВ}_\textsl{пасп}$}$,
%$\text{$\textsl{КРВ}_\textsl{пасп\_Iном}$}$,
%$\text{$\textsl{КРВ}_\textsl{пасп\_Iном\_откл}$}$,
%$I_{\textsl{ном\_В\_пасп}}$,
%$I_{\textsl{ном\_откл\_В\_пасп}}$
%\end{minipage}}уставки функции.
\end{eqexpl}

\vspace{3mm}

\item
Если выполняется неравенство $I_{\textsl{макс\_выч.}} > I_{\textsl{ном\_откл\_В\_пасп}}$, то выключатель считается выработавшим свой ресурс, а «КРВ» приравнивается к нулю. При помощи сигнала «Сброс КРВ» счетчики «КРВ» и «МРВ» сбрасываются на начальные значения. В функции предусмотрено поле, в котором отображаются текущие значения «КРВ» и «МРВ». Текущие значения остаточного ресурса выключателя можно узнать через соответствующее меню ИЧМ и в ПО ЮНИТ-Сервис.

\item
Ввод функции в работу на этапе параметрирования устройства осуществляется программным переключателем SGF1 «Ввод\_функции». Информация о введенной функции отображается выходным сигналом «КРВ: Ввод».

\item
Алгоритм работы функции выполнен в соответствии с рисунком \ref{krv:img1}. В таблице \ref{krv:tbl1} приведены параметры, необходимые для настройки функции.

\vspace{3mm}
\begin{figure}[H]
\includegraphics[width=\textwidth,height=\textheight,keepaspectratio]{\fbpath/30. КРВ/_latex/img/1.pdf}
\captionof{figure}{Функциональная схема логики КРВ}
\label{krv:img1}
\end{figure}

\end{enumerate}
\FloatBarrier % Все float'ы ДО этой строки будут размещены здесь.
